% ============================================================
% Chapter 11: Detection, Threat Hunting, Incident Response and Lab Operations
% Language: Greek — edit freely; regenerate from src/content if preferred.
% ============================================================
\chapter{Ανίχνευση, Προληπτική Αναζήτηση Απειλών, Απόκριση σε Περιστατικά και Εργαστηριακές Λειτουργίες}\label{ch:11}
\chaptermeta{IDS/IPS · Υπογραφές και ανίχνευση ανωμαλιών · Τεχνικές αποφυγής · Προληπτική αναζήτηση απειλών · YARA και Sigma · Απομονωμένα εργαστήρια · Απόκριση σε περιστατικά κατά NIST · Λειτουργίες SOC}{Επίπεδο: Μεσαίο–Προχωρημένο επίπεδο · Λειτουργίες}{Χρόνος μελέτης: 10–12 ώρες}

Η πρόληψη αργά ή γρήγορα αποτυγχάνει. Το Κεφάλαιο 1 εισήγαγε την άμυνα σε βάθος ακριβώς επειδή κάθε προληπτικό μέτρο έχει κενά, και τα προηγούμενα κεφάλαια έδειξαν πώς τα εντοπίζουν οι ικανοί αντίπαλοι. Το παρόν κεφάλαιο αφορά, επομένως, όσα ακολουθούν: πώς ένας οργανισμός \emph{ανιχνεύει} κακόβουλη δραστηριότητα, \emph{αναζητά} προληπτικά απειλές που διέφυγαν από τους συναγερμούς του και \emph{αποκρίνεται} με πειθαρχημένο τρόπο που περιορίζει τη ζημιά και διαφυλάσσει τα αποδεικτικά στοιχεία.

Ξεκινάμε με τα συστήματα ανίχνευσης και πρόληψης εισβολών, τα δύο κύρια παραδείγματα επιθεώρησης και τις τεχνικές με τις οποίες οι επιτιθέμενοι τα αποφεύγουν. Στη συνέχεια στρεφόμαστε στην προληπτική αναζήτηση απειλών, στις γλώσσες κανόνων YARA και Sigma και στον σχεδιασμό απομονωμένων εργαστηρίων, όπου το κακόβουλο λογισμικό μπορεί να μελετηθεί με ασφάλεια. Το κεφάλαιο ολοκληρώνεται με τον κύκλο απόκρισης σε περιστατικά όπως τον ορίζει το NIST και με την οργάνωση ενός Κέντρου Επιχειρήσεων Ασφάλειας (SOC).

\begin{outcomesbox}{Μαθησιακά αποτελέσματα}
Μετά τη μελέτη του κεφαλαίου θα πρέπει να μπορείτε:
\begin{itemize}
  \item Να συγκρίνετε τα δικτυακά και τα τοπικά IDS/IPS, καθώς και την επιθεώρηση βάσει υπογραφών και βάσει ανωμαλιών.
  \item Να εξηγείτε συνήθεις τεχνικές αποφυγής των IDS και τον ρόλο της κανονικοποίησης της κίνησης.
  \item Να σχεδιάζετε μια προληπτική αναζήτηση απειλών βασισμένη σε υποθέσεις, αξιοποιώντας τηλεμετρία τερματικών και μνήμης.
  \item Να συντάσσετε βασικούς κανόνες YARA και Sigma και να διαχειρίζεστε τις ανιχνεύσεις ως κώδικα.
  \item Να εφαρμόζετε τον κύκλο απόκρισης σε περιστατικά του NIST και να περιγράφετε τις βαθμίδες ενός SOC και τη ροή ενός SIEM.
\end{itemize}
\end{outcomesbox}

\section{Αρχιτεκτονικές IDS και IPS και παραδείγματα επιθεώρησης}\label{sec:11.1}
Ένα \textbf{σύστημα ανίχνευσης εισβολών} (Intrusion Detection System, IDS) παρακολουθεί τη δραστηριότητα και εκδίδει ειδοποιήσεις· ένα \textbf{σύστημα πρόληψης εισβολών} (Intrusion Prevention System, IPS) βρίσκεται στη διαδρομή της κίνησης και μπορεί να αποκλείει αυτόματα την κακόβουλη δραστηριότητα. Ένα \textbf{δικτυακό IDS} (NIDS), όπως τα Suricata, Snort ή Zeek, αναλύει αντίγραφα της δικτυακής κίνησης από θύρα κατοπτρισμού ενός μεταγωγέα ή από δικτυακό διακλαδωτή (tap) και βλέπει την επικοινωνία μεταξύ πολλών υπολογιστών. Ένα \textbf{τοπικό IDS} (HIDS), όπως τα Wazuh ή OSSEC, εκτελείται σε μεμονωμένα συστήματα και παρατηρεί αλλαγές αρχείων, αρχεία καταγραφής και διεργασίες —συμπεριλαμβανομένης δραστηριότητας μέσα σε κρυπτογραφημένες συνδέσεις που ένας δικτυακός αισθητήρας δεν μπορεί να διαβάσει. Οι δύο οπτικές είναι συμπληρωματικές.

Η λογική της ανίχνευσης ακολουθεί δύο παραδείγματα. Η ανίχνευση \textbf{βάσει υπογραφών} αναζητά γνωστά μοτίβα, όπως μια συγκεκριμένη συμβολοσειρά εκμετάλλευσης· είναι ακριβής και εύκολα εξηγήσιμη, αλλά τυφλή απέναντι σε νέες επιθέσεις. Η ανίχνευση \textbf{βάσει ανωμαλιών ή συμπεριφοράς} μοντελοποιεί το φυσιολογικό και επισημαίνει αποκλίσεις, όπως έναν σταθμό εργασίας που ξαφνικά μεταφέρει gigabytes τη νύχτα· μπορεί να εντοπίσει άγνωστες επιθέσεις, αλλά παράγει περισσότερα ψευδώς θετικά αποτελέσματα και απαιτεί προσεκτικά καθορισμένη συμπεριφορά αναφοράς. Τα ώριμα προγράμματα συνδυάζουν και τα δύο και τα βελτιστοποιούν συνεχώς, επειδή μια ειδοποίηση που οι αναλυτές έχουν μάθει να αγνοούν είναι χειρότερη από την απουσία ειδοποίησης.

\section{Τεχνικές αποφυγής και αυτοματοποιημένη αντιμετώπιση}\label{sec:11.2}
Οι επιτιθέμενοι προσπαθούν να κάνουν το IDS να «βλέπει» κάτι διαφορετικό από αυτό που λαμβάνει πράγματι ο στόχος. Ο \textbf{κατακερματισμός} διασπά ένα κακόβουλο φορτίο σε πολλά μικρά πακέτα ή σε επικαλυπτόμενα τμήματα· η \textbf{κωδικοποίηση} και η συσκότιση μεταμφιέζουν γνωστές συμβολοσειρές· η χειραγώγηση του χρονισμού ή των τιμών \emph{TTL} μπορεί να κάνει τον αισθητήρα και τον προορισμό να ανασυνθέτουν μια ροή με διαφορετικό τρόπο· και, όλο και συχνότερα, οι επιθέσεις απλώς κρύβονται μέσα σε \textbf{κρυπτογραφημένη} κίνηση TLS. Τα αντίμετρα περιλαμβάνουν την \textbf{κανονικοποίηση της κίνησης}, η οποία ανασυνθέτει τις ροές με τον ίδιο τρόπο όπως ο τελικός υπολογιστής, την επιθεώρηση TLS σε ελεγχόμενα σημεία όπου αυτό είναι νόμιμο και αναλογικό, καθώς και τη μεγαλύτερη αξιοποίηση της τοπικής τηλεμετρίας και των μεταδεδομένων, όπως ο χρονισμός των συνδέσεων.

Όταν ένα IPS ή ένα αυτοματοποιημένο σχέδιο ενεργειών αποκρίνεται, μπορεί να \textbf{απορρίψει} τα επίμαχα πακέτα, να \textbf{επαναφέρει} τη σύνδεση TCP, να \textbf{αποκλείσει} τη διεύθυνση προέλευσης για ένα διάστημα, να \textbf{αναδιαμορφώσει} ένα τείχος προστασίας ή να \textbf{απομονώσει} ένα τερματικό από το δίκτυο. Ο αυτοματισμός μειώνει δραματικά τον χρόνο απόκρισης, όμως ένας λανθασμένα ρυθμισμένος κανόνας μπορεί επίσης να διακόψει νόμιμη επιχειρησιακή δραστηριότητα. Γι' αυτό ο αυτόματος αποκλεισμός επιφυλάσσεται για ανιχνεύσεις υψηλής βεβαιότητας, ενώ οι ειδοποιήσεις χαμηλότερης βεβαιότητας διοχετεύονται στους αναλυτές.

\section{Προληπτική αναζήτηση απειλών και τηλεμετρία τερματικών}\label{sec:11.3}
Η \textbf{προληπτική αναζήτηση απειλών} (threat hunting) ξεκινά από την παραδοχή ότι ένας αντίπαλος μπορεί να βρίσκεται ήδη στο περιβάλλον χωρίς να έχει ενεργοποιήσει κανέναν συναγερμό. Η αναζήτηση καθοδηγείται από μια \textbf{υπόθεση}, που συνήθως προέρχεται από πληροφορίες απειλών ή από το ATT\&CK —για παράδειγμα «ένας επιτιθέμενος χρησιμοποιεί προγραμματισμένες εργασίες για μόνιμη παρουσία στους διακομιστές μας». Ο αναλυτής υποβάλλει στη συνέχεια ερωτήματα στη διαθέσιμη τηλεμετρία, διερευνά τις ανωμαλίες και καταλήγει σε συμπέρασμα. Όποιο κι αν είναι το αποτέλεσμα, μια καλή αναζήτηση ολοκληρώνεται με τη μετατροπή της λογικής της σε αυτοματοποιημένη ανίχνευση, ώστε η ίδια συμπεριφορά να εντοπίζεται στο μέλλον.

Η αναζήτηση απειλών εξαρτάται από πλούσια \textbf{τηλεμετρία τερματικών}. Εργαλεία όπως το Sysmon της Microsoft καταγράφουν τη δημιουργία διεργασιών με την πλήρη γραμμή εντολών, τις δικτυακές συνδέσεις, τη φόρτωση οδηγών και τις αλλαγές στο μητρώο· οι πλατφόρμες EDR συλλέγουν παρόμοια δεδομένα σε μεγάλη κλίμακα· και στο Linux, το \texttt{auditd} και οι αισθητήρες που βασίζονται στο eBPF παρέχουν ισοδύναμη ορατότητα. Η \textbf{ανάλυση μνήμης} με πλαίσια όπως το Volatility μπορεί να αποκαλύψει εγχυμένο κώδικα, κρυφές διεργασίες και αποκρυπτογραφημένες ρυθμίσεις κακόβουλου λογισμικού που δεν γράφτηκαν ποτέ στον δίσκο —ικανότητα απαραίτητη απέναντι στις τεχνικές χωρίς αρχεία που περιγράφηκαν στο Κεφάλαιο 4.

\section{YARA, Sigma και ανιχνεύσεις ως κώδικας}\label{sec:11.4}
Η \textbf{YARA} είναι μια γλώσσα αντιστοίχισης προτύπων για την ταξινόμηση αρχείων και περιοχών μνήμης. Ένας κανόνας έχει τρία μέρη: \emph{meta} (περιγραφικές πληροφορίες), \emph{strings} (πρότυπα κειμένου, δεκαεξαδικά ή κανονικές εκφράσεις) και μια \emph{condition} που τα συνδυάζει, για παράδειγμα \texttt{uint16(0) == 0x5A4D and 2 of (\$s*)}, που σημαίνει «εκτελέσιμο Windows που περιέχει τουλάχιστον δύο από τις δηλωμένες συμβολοσειρές». Οι καλοί κανόνες στοχεύουν χαρακτηριστικά που ο δημιουργός του κακόβουλου λογισμικού δεν μπορεί να αλλάξει εύκολα —μοναδικές ακολουθίες κώδικα, δομές ρυθμίσεων— και όχι μία μόνο συμβολοσειρά, και δοκιμάζονται σε συλλογή νόμιμων αρχείων για την αποφυγή ψευδώς θετικών αποτελεσμάτων.

Ενώ η YARA επιθεωρεί αρχεία, η \textbf{Sigma} περιγράφει ύποπτα μοτίβα σε \emph{συμβάντα αρχείων καταγραφής}, σε μορφή ανεξάρτητη από προμηθευτή, η οποία μπορεί να μετατραπεί σε ερωτήματα για διαφορετικές πλατφόρμες SIEM. Η προσέγγιση των \textbf{ανιχνεύσεων ως κώδικα} (detection-as-code) εφαρμόζει και στις δύο την πειθαρχία της μηχανικής λογισμικού: οι κανόνες αποθηκεύονται σε σύστημα ελέγχου εκδόσεων, ελέγχονται από ομοτίμους, δοκιμάζονται αυτόματα σε δείγματα δεδομένων, αντιστοιχίζονται σε τεχνικές του ATT\&CK και εγκαθίστανται μέσω μιας ροής ανάπτυξης. Οι πλατφόρμες \textbf{SOAR} (ενορχήστρωση, αυτοματοποίηση και απόκριση ασφάλειας) εκτελούν στη συνέχεια σχέδια ενεργειών όταν ενεργοποιείται μια ανίχνευση —εμπλουτίζοντας την ειδοποίηση, ανοίγοντας ένα δελτίο και, αν η βεβαιότητα είναι υψηλή, απομονώνοντας τον υπολογιστή.

\section{Σχεδιασμός απομονωμένων εργαστηρίων ανάλυσης}\label{sec:11.5}
Η μηχανική ανιχνεύσεων και η ανάλυση κακόβουλου λογισμικού απαιτούν έναν χώρο όπου ο εχθρικός κώδικας μπορεί να εκτελεστεί χωρίς να θέτει σε κίνδυνο πραγματικά συστήματα. Ένα ασφαλές εργαστήριο βασίζεται στην \textbf{εικονικοποίηση}: τα μηχανήματα ανάλυσης λειτουργούν ως εικονικές μηχανές σε αποκλειστικό φυσικό υπολογιστή, με \textbf{στιγμιότυπα} (snapshots) που επιτρέπουν την επαναφορά τους σε καθαρή κατάσταση μετά από κάθε πείραμα. Η δικτύωση ορίζεται ως \emph{host-only} ή ως εσωτερικό δίκτυο χωρίς διαδρομή προς το διαδίκτυο ή το περιβάλλον παραγωγής· όπου πρέπει να παρατηρηθεί η δικτυακή συμπεριφορά, προσομοιωμένες υπηρεσίες όπως τα INetSim ή FakeNet απαντούν στα αιτήματα του κακόβουλου λογισμικού.

Η απομόνωση πρέπει να είναι και διαδικαστική. Οι κοινόχρηστοι φάκελοι και ο συγχρονισμός του προχείρου μεταξύ φυσικού και εικονικού μηχανήματος πρέπει να απενεργοποιούνται, τα δείγματα να αποθηκεύονται σε αρχεία συμπίεσης με κωδικό και με σαφή σήμανση, ενώ οι αναλυτές πρέπει να θυμούνται ότι ορισμένα κακόβουλα προγράμματα ανιχνεύουν τις εικονικές μηχανές και αλλάζουν συμπεριφορά. Οι προφυλάξεις αυτές εφαρμόζουν την ίδια λογική περιορισμού που διέπει την απόκριση σε περιστατικά.

\section{Ο κύκλος απόκρισης σε περιστατικά (NIST SP 800-61)}\label{sec:11.6}
Η Ειδική Έκδοση 800-61 του NIST περιγράφει την απόκριση σε περιστατικά ως κύκλο τεσσάρων φάσεων. Πρώτη έρχεται η \textbf{προετοιμασία}, η οποία καθορίζει όλα τα υπόλοιπα: εγκεκριμένο σχέδιο, καθορισμένοι ρόλοι, κατάλογοι επαφών, πρότυπα επικοινωνίας, εργαλεία, καταγραφή και τακτικές ασκήσεις. Η \textbf{ανίχνευση και ανάλυση} επιβεβαιώνει ότι σημειώθηκε περιστατικό, προσδιορίζει το εύρος και τη σοβαρότητά του και τεκμηριώνει κάθε βήμα. Ο \textbf{περιορισμός, η εξάλειψη και η ανάκαμψη} περιορίζουν τη ζημιά —βραχυπρόθεσμα με την απομόνωση συστημάτων και την απενεργοποίηση παραβιασμένων λογαριασμών, μακροπρόθεσμα με την ανακατασκευή τους— αφαιρούν τα ερείσματα του επιτιθέμενου και αποκαθιστούν την κανονική λειτουργία από αξιόπιστες πηγές. Η \textbf{δραστηριότητα μετά το περιστατικό} περιλαμβάνει μια ανασκόπηση διδαγμάτων χωρίς απόδοση ευθυνών μέσα σε λίγες εβδομάδες και επιστρέφει τις βελτιώσεις στη φάση της προετοιμασίας.

Δύο εντάσεις διατρέχουν κάθε περιστατικό. Η πρώτη είναι μεταξύ \textbf{ταχύτητας και αποδεικτικών στοιχείων}: η απενεργοποίηση ενός μηχανήματος σταματά την επίθεση, αλλά καταστρέφει τα περιεχόμενα της μνήμης· γι' αυτό τα αποδεικτικά στοιχεία πρέπει, όπου είναι δυνατόν, να συλλέγονται πριν από τις ενέργειες που διαταράσσουν το σύστημα (Κεφάλαιο 12). Η δεύτερη είναι μεταξύ \textbf{περιορισμού και συλλογής πληροφοριών}: η πρόωρη απομάκρυνση του επιτιθέμενου, πριν γίνει γνωστό ολόκληρο το αποτύπωμά του, μπορεί απλώς να τον προειδοποιήσει και να τον ωθήσει να χρησιμοποιήσει άλλη κερκόπορτα. Η αναθεωρημένη καθοδήγηση (Rev. 3, 2025) ευθυγραμμίζει αυτές τις δραστηριότητες με τις λειτουργίες του NIST Cybersecurity Framework 2.0.

\begin{notebox}{Νομικές και κανονιστικές προθεσμίες}
Πολλά περιστατικά ενεργοποιούν υποχρεωτικές γνωστοποιήσεις —για παράδειγμα, εντός 72 ωρών στην εποπτική αρχή βάσει του ΓΚΠΔ για παραβιάσεις προσωπικών δεδομένων και έγκαιρη προειδοποίηση εντός 24 ωρών βάσει της NIS2 για σημαντικά περιστατικά (Κεφάλαιο 13). Το σχέδιο απόκρισης πρέπει να προσδιορίζει ποιος αποφασίζει και ποιος προβαίνει στη γνωστοποίηση.
\end{notebox}

\section{Λειτουργίες SOC: βαθμίδες, η ροή του SIEM και σχεδιασμός σεναρίων χρήσης}\label{sec:11.7}
Ένα \textbf{Κέντρο Επιχειρήσεων Ασφάλειας} (Security Operations Centre, SOC) είναι η ομάδα και η υποδομή που παρακολουθούν συνεχώς έναν οργανισμό. Η εργασία κατανέμεται παραδοσιακά σε βαθμίδες: οι αναλυτές της \textbf{Βαθμίδας 1} διαλογούν τις εισερχόμενες ειδοποιήσεις, η \textbf{Βαθμίδα 2} διερευνά σε βάθος τα επιβεβαιωμένα περιστατικά και οι ειδικοί της \textbf{Βαθμίδας 3} αναζητούν απειλές, αναλύουν κακόβουλο λογισμικό και σχεδιάζουν ανιχνεύσεις. Το κεντρικό εργαλείο είναι η πλατφόρμα \textbf{SIEM} (διαχείριση πληροφοριών και συμβάντων ασφάλειας), της οποίας η ροή \emph{συλλέγει} αρχεία καταγραφής από πολλές πηγές, τα \emph{κανονικοποιεί} σε ένα κοινό σχήμα, τα \emph{εμπλουτίζει} με πλαίσιο, όπως οι υπεύθυνοι των στοιχείων και οι πληροφορίες απειλών, \emph{συσχετίζει} συμβάντα σε ειδοποιήσεις και τα \emph{αποθηκεύει} για διερεύνηση και κανονιστική συμμόρφωση.

Ένα SIEM είναι τόσο καλό όσο τα \textbf{σενάρια χρήσης} του —οι συγκεκριμένες απειλές που έχει ρυθμιστεί να ανιχνεύει. Κάθε σενάριο χρήσης πρέπει να δηλώνει την απειλή που αντιμετωπίζει (ιδανικά ως τεχνική του ATT\&CK), τις πηγές δεδομένων που απαιτεί, τη λογική ανίχνευσης, το αναμενόμενο ποσοστό ψευδώς θετικών και τη διαδικασία απόκρισης. Η απόδοση ενός SOC μετράται συνήθως μέσω του \textbf{μέσου χρόνου ανίχνευσης (MTTD)} και του \textbf{μέσου χρόνου απόκρισης (MTTR)}, σε συνδυασμό με το ποσοστό των ειδοποιήσεων που αποδεικνύονται αξιοποιήσιμες —βασικό δείκτη του φόρτου εργασίας και της εξουθένωσης των αναλυτών.

\section*{Βασικοί όροι}
\begin{description}[style=nextline,leftmargin=1.2em]
  \item[NIDS / HIDS] Δικτυακά και τοπικά συστήματα ανίχνευσης εισβολών.
  \item[Κανονικοποίηση κίνησης] Ανασύνθεση της κίνησης όπως θα την ανασυνέθετε ο προορισμός, για την ακύρωση τεχνικών αποφυγής.
  \item[Προληπτική αναζήτηση απειλών] Προληπτική αναζήτηση αντιπάλων που διέφυγαν της ανίχνευσης, καθοδηγούμενη από υποθέσεις.
  \item[YARA / Sigma] Γλώσσες κανόνων για την αντιστοίχιση προτύπων σε αρχεία (YARA) και σε συμβάντα καταγραφής (Sigma).
  \item[Περιορισμός] Ενέργειες απόκρισης που περιορίζουν τη διάδοση και την επίπτωση μιας επίθεσης.
  \item[SIEM] Πλατφόρμα που συλλέγει, κανονικοποιεί, συσχετίζει και αποθηκεύει συμβάντα ασφάλειας.
\end{description}

\section*{Σύνοψη κεφαλαίου}
\begin{itemize}
  \item Οι δικτυακοί και οι τοπικοί αισθητήρες παρέχουν συμπληρωματική ορατότητα· η ανίχνευση βάσει υπογραφών και βάσει ανωμαλιών έχει η καθεμιά πλεονεκτήματα και αδυναμίες.
  \item Οι επιτιθέμενοι αποφεύγουν την επιθεώρηση μέσω κατακερματισμού, κωδικοποίησης και κρυπτογράφησης· η κανονικοποίηση και η τοπική τηλεμετρία τούς αντιμετωπίζουν.
  \item Η προληπτική αναζήτηση απειλών καθοδηγείται από υποθέσεις και πρέπει να καταλήγει στη δημιουργία νέων αυτοματοποιημένων ανιχνεύσεων.
  \item Οι κανόνες YARA και Sigma, όταν διαχειρίζονται ως κώδικας και συνδέονται με SOAR, μετατρέπουν τη γνώση σε επαναλήψιμη ανίχνευση και απόκριση.
  \item Η απόκριση σε περιστατικά ακολουθεί τον κύκλο του NIST, εξισορροπώντας την ταχύτητα με τη διαφύλαξη των αποδεικτικών στοιχείων, και το SOC την υλοποιεί μέσω βαθμίδων, ενός SIEM και μετρήσιμων σεναρίων χρήσης.
\end{itemize}

\section*{Ερωτήσεις ανασκόπησης}
\begin{enumerate}
  \item Γιατί ένας τοπικός αισθητήρας μπορεί να ανιχνεύσει μια επίθεση μέσα σε κίνηση HTTPS που διαφεύγει από ένα δικτυακό IDS;
  \item Διατυπώστε μια υπόθεση αναζήτησης για την υποκλοπή διαπιστευτηρίων από τη διεργασία LSASS και αναφέρετε την τηλεμετρία που θα χρειαζόσασταν.
  \item Γιατί ένας κανόνας YARA δεν πρέπει να βασίζεται σε μία μόνο συμβολοσειρά, όπως ένα όνομα τομέα C2;
  \item Κατά τη διάρκεια ενός περιστατικού, ένας προϊστάμενος θέλει να απενεργοποιήσει αμέσως τον επηρεαζόμενο διακομιστή. Συζητήστε τα αντισταθμιστικά οφέλη και κόστη.
\end{enumerate}
